\chapter{Connecting a Shake}\label{ch:connect}

\section{Cabling and power}
\begin{enumerate}
\item Place the Shake on a solid, level floor, ideally on the ground floor, away from doors, fans and footsteps.
\item Connect its Ethernet port directly to your laptop's network port.
\item Power the Shake with its own supply and wait \textbf{three to five minutes} for it to start.
\end{enumerate}
Without a router (no DHCP server) both ends give themselves a \emph{link-local} address in the range \code{169.254.x.x}. That
is normal and is what RS3D expects.

\section{Give the laptop a 169.254.x.x address}
The laptop needs an address in the same range on the wired port. Many systems do this by themselves when the cable is plugged in.
If they do not, the \ui{Status} page shows \ui{Laptop IP} in red with a \ui{Set IP} button
(Chapter~\ref{ch:status}). Pressing it shows the exact commands, asks your system for permission and applies them.

\section{Find the Shake}
Press \ui{Find Shake\ldots} in the sidebar. RS3D looks for Shakes on the cable (it tries the names \code{rs.local},
\code{rs-2.local} and so on, and the addresses your computer has recently seen) and confirms each candidate by asking it for
its station name. Every Shake it finds is listed with its station code and address; press \ui{Add} to add it.

If nothing is found, check the cable, wait a little longer for the Shake to start, and see Chapter~\ref{ch:trouble}; you can always add
the Shake by hand.

\section{Add a Shake by hand}
Open \ui{Settings}. The \ui{Device} card has these fields.

\begin{figure}[tbp]
\centering
\annotshot[0.82\linewidth]{settings}{\calloutsSettings}
\caption{The \ui{Settings} page. (1)--(6) The device fields, described below. (7) \ui{Save}. (8) \ui{Test connection}. (9) Time zone of the displayed times. (10) The data folder.}
\end{figure}

\begin{enumerate}
\item \ui{Name}: what you want to call it in the sidebar (default: the station code).
\item \ui{IP / host name}: the Shake's address, for example \code{169.254.221.62} or \code{rs.local}.
\item \ui{SeedLink port}: the live-data port; 18000 unless you changed it on the Shake.
\item \ui{SSH user}: the account RS3D uses to read the Shake's clock, storage and archive; \code{myshake} on a standard Shake.
\item \ui{SSH password}: that account's password. It is stored in your operating system's key store (not in a file) and never
  leaves your computer.
\item \ui{SSH key file (optional)}: a private key, if you prefer key login to a password.
\end{enumerate}
Press \ui{Test connection}; \ui{Connection OK} means the login works. Then press \ui{Save}.

\begin{rs3dwarn}[Change the default password]
A new Raspberry Shake uses a published default password. The \ui{Password} check on the Status page warns while it is unchanged. To
change it, log in to the Shake over SSH and run \code{passwd}, then enter the new password in \ui{Settings}. RS3D only \emph{reads}
from the Shake and never changes anything on it.
\end{rs3dwarn}

\section{More than one Shake}
Every Shake you add appears in the sidebar under \ui{Devices}; click a name to switch. \ui{New device} on the \ui{Settings} page
starts an empty form, and \ui{Remove} deletes the selected device (its downloaded data stay on disk).
